\section{Partial Derivatives}

For $f : \R^n \to \R$ (often $n = 2$ or $3$).

\subsection{Limits and continuity}
$\displaystyle \lim_{(x,y)\to(a,b)} f(x,y) = L$ requires the same value along \emph{every}
path of approach. To show a limit does not exist: find two paths giving different limits.

\begin{example}[Path-dependence]
$f(x,y) = \tfrac{xy}{x^2+y^2}$ at $(0,0)$: along $y = 0$, $f \equiv 0$; along $y = x$, $f = 1/2$.
Limit does not exist.
\end{example}

\subsection{Partial derivatives}
\[
\frac{\partial f}{\partial x}(a, b) = \lim_{h \to 0} \frac{f(a + h, b) - f(a, b)}{h}.
\]
Compute by treating other variables as constants. Higher-order partials: $f_{xy} = (f_x)_y$.

\begin{theorem}[Clairaut]
If $f_{xy}$ and $f_{yx}$ are continuous on an open set containing $(a,b)$, then
$f_{xy}(a,b) = f_{yx}(a,b)$.
\end{theorem}

\subsection{Differentiability and tangent plane}
$f$ is \emph{differentiable} at $(a,b)$ if there exists a linear map $L$ such that
\[ f(a + h, b + k) = f(a,b) + L(h,k) + o(\sqrt{h^2 + k^2}). \]
If $f_x, f_y$ exist and are continuous in a neighborhood of $(a,b)$, then $f$ is differentiable there,
and the tangent plane to $z = f(x,y)$ is
\[ z = f(a,b) + f_x(a,b) (x - a) + f_y(a,b) (y - b). \]

\subsection{Total differential and linear approximation}
\[ df = f_x \dx + f_y \dy + f_z \dz, \qquad f(\vc x + d\vc x) \approx f(\vc x) + \grad f(\vc x) \dotp d\vc x. \]

\subsection{Chain rule}
\textbf{One independent variable:} $z = f(x(t), y(t))$,
\[ \frac{dz}{dt} = \frac{\partial f}{\partial x}\frac{dx}{dt} + \frac{\partial f}{\partial y}\frac{dy}{dt}. \]

\textbf{Several independent variables:} $z = f(x(s,t), y(s,t))$,
\[ \frac{\partial z}{\partial s} = f_x \frac{\partial x}{\partial s} + f_y \frac{\partial y}{\partial s}. \]

\textbf{General matrix form:} for $f: \R^n \to \R^m$ and $g: \R^k \to \R^n$,
\[ J_{f \circ g}(\vc x) = J_f(g(\vc x)) \, J_g(\vc x). \]

\subsection{Implicit differentiation}
If $F(x, y, z) = 0$ defines $z$ implicitly as a function of $x, y$:
\[ \frac{\partial z}{\partial x} = -\frac{F_x}{F_z}, \qquad \frac{\partial z}{\partial y} = -\frac{F_y}{F_z}, \qquad (F_z \neq 0). \]

\begin{pitfallbox}
Existence of $f_x$ and $f_y$ at a point does \emph{not} imply $f$ is even continuous there.
Continuous partials in a neighborhood is the safe sufficient condition for differentiability.
\end{pitfallbox}
